\documentclass{standalone}
\usepackage{circuitikz}
\usetikzlibrary{calc}
\definecolor{circuitnotered}{HTML}{E5534B}
\definecolor{circuitnoteblue}{HTML}{4C8EDA}
\definecolor{circuitnotegreen}{HTML}{2EA043}
\definecolor{circuitnoteorange}{HTML}{D29922}
\definecolor{circuitnotemark}{HTML}{FE00FE}
\begin{document}
\begin{circuitikz}[american, line width=0.8pt]
\ctikzset{bipoles/length=1.2cm}
\foreach \x in {0,2,4,6,8,10} {\foreach \y in {0,-2,-4} {\fill[gray, opacity=0.35] (\x,\y) circle (1.1pt);}}
\node[gray, font=\scriptsize] at (0,0.84) {1};
\node[gray, font=\scriptsize] at (2,0.84) {2};
\node[gray, font=\scriptsize] at (4,0.84) {3};
\node[gray, font=\scriptsize] at (6,0.84) {4};
\node[gray, font=\scriptsize] at (8,0.84) {5};
\node[gray, font=\scriptsize] at (10,0.84) {6};
\node[gray, font=\scriptsize, anchor=east] at (-0.84,0) {a};
\node[gray, font=\scriptsize, anchor=east] at (-0.84,-2) {b};
\node[gray, font=\scriptsize, anchor=east] at (-0.84,-4) {c};
\coordinate (x1y1) at (0,0);
\coordinate (x3y1) at (4,0);
\coordinate (x4y1) at (6,0);
\coordinate (x6y1) at (10,0);
\coordinate (x1y3) at (0,-4);
\coordinate (x3y3) at (4,-4);
\coordinate (x4y3) at (6,-4);
\coordinate (x6y3) at (10,-4);
\draw (x1y1) to[R, l_=$R_{1}$] (x3y1); % line 3
\draw (x4y1) to[R, l_=$R_{2}$] (x6y1); % line 4
\draw (x1y3) to[R, l_=$R_{3}$] (x3y3); % line 5
\draw (x4y3) to[R, l_=$R_{4}$] (x6y3); % line 6
\draw[circuitnotered] (2,0) circle (0.9); % line 8
\draw[circuitnoteblue] (8,0) circle (0.9); % line 9
\draw[circuitnotegreen] (2,-4) circle (0.9); % line 10
\draw[circuitnoteorange] (8,-4) circle (0.9); % line 11
\path (1.619,-1.993) rectangle (2.381,-0.807); % line 12
\node[anchor=west, inner sep=0, circuitnotemark, font=\large] at (2,-1.4) {X}; % line 12
\path (7.492,-1.993) rectangle (8.508,-0.807); % line 13
\node[anchor=west, inner sep=0, circuitnotemark, font=\large] at (8,-1.4) {X}; % line 13
\path (0,-2.593) rectangle (4.572,-1.407); % line 14
\node[anchor=west, inner sep=0, circuitnotemark, font=\large] at (0,-2) {X}; % line 14
\path (6,-2.593) rectangle (10.572,-1.407); % line 15
\node[anchor=west, inner sep=0, circuitnotemark, font=\large] at (6,-2) {X}; % line 15
\path (1.365,-5.993) rectangle (2.635,-4.807); % line 16
\node[anchor=west, inner sep=0, circuitnotemark, font=\large] at (2,-5.4) {X}; % line 16
\path (7.238,-5.993) rectangle (8.762,-4.807); % line 17
\node[anchor=west, inner sep=0, circuitnotemark, font=\large] at (8,-5.4) {X}; % line 17
\path (0,-6.593) rectangle (4.572,-5.407); % line 18
\node[anchor=west, inner sep=0, circuitnotemark, font=\large] at (0,-6) {X}; % line 18
\path (6,-6.593) rectangle (10.572,-5.407); % line 19
\node[anchor=west, inner sep=0, circuitnotemark, font=\large] at (6,-6) {X}; % line 19
\node[anchor=south west, inner sep=2pt, yshift=4pt, circuitnotemark, font=\LARGE\bfseries] (circuittitle) at (current bounding box.north west) {X};
\path (circuittitle.south west) rectangle ++(4.938,0.853);
\node[anchor=north east, inner sep=2pt, circuitnotemark, font=\large] (circuitstamp) at ([yshift=-0.593cm]current bounding box.south east) {X};
\path (circuitstamp.north east) rectangle ++(-5.08,-0.593);
\end{circuitikz}
\end{document}
